\documentclass[11pt]{article}
\usepackage[T1]{fontenc}
\usepackage{lmodern}
\usepackage[margin=1in]{geometry}
\usepackage{mathtools,amssymb,amsthm}
\usepackage{microtype}
\usepackage[hidelinks]{hyperref}
\hypersetup{
 pdftitle={Higher-index failure of ultra-log-concavity for Chen's recursion polynomials},
 pdfauthor={Xiaojian Zeng},
 pdfsubject={A proof of arXiv:2503.16296v2, Conjecture 3.9},
 pdfkeywords={ultra-log-concavity, recursion polynomials, signed moments}}
\newtheorem{theorem}{Theorem}
\newtheorem{lemma}[theorem]{Lemma}
\newcommand{\dd}{\,\mathrm{d}}
\newcommand{\ULC}{\mathrm{ULC}}
\title{Higher-index failure of ultra-log-concavity\newline
for Chen's recursion polynomials}
\author{Xiaojian Zeng\thanks{PBC School of Finance, Tsinghua University,
100190, China. Email: \href{mailto:xiaojianzeng@foxmail.com}{xiaojianzeng@foxmail.com}.}}
\date{}
\begin{document}
\maketitle

\begin{abstract}
For the degree-$(m-1)$ polynomials
$f_m(s)=((s+1)^m-(-1)^m)/(s+2)$ and
$h_m(s)=((s+1)^m+(-1)^m(s+1))/(s+2)$, we prove that every
ordinary ultra-log-concavity inequality at an index
$4\le j\le m-2$ fails strictly. This proves Conjecture~3.9 of
Chen's arXiv version~2 on log-concavity for banana graphs and clasped
necklaces. The normalized coefficient difference reduces exactly to a
two-by-two moment determinant on $[-1,1]$. For even parity the determinant
is positive by Cauchy--Schwarz. For odd parity, where the density is signed,
an explicit base case and an identity involving an auxiliary positive
measure give a positivity-preserving induction. An endpoint difference
equal to $-m$ also recovers the known failure of ultra-log-concavity for
every $m\ge4$, including the two cases with an empty higher-index range.
\end{abstract}

\section{Statement and context}

Aluffi, Marcolli and Qaisar developed recursive formulas for the
Grothendieck classes of melonic graphs and proved positivity of their
coefficients in the variable given by the class of $\mathcal M_{0,4}$
\cite[Proposition~3.2 and Corollary~3.4]{AMQ2023}.
Two of the recursion polynomials studied by Chen are
\begin{equation}\label{eq:families}
 f_m(s)=\frac{(s+1)^m-(-1)^m}{s+2},\qquad
 h_m(s)=\frac{(s+1)^m+(-1)^m(s+1)}{s+2}.
\end{equation}
Chen proves ordinary coefficient log-concavity for these families and
investigates the stronger ultra-log-concavity inequalities
\cite[Theorem~1.1 and Section~3]{Chen2025}.

For a polynomial $P(s)=\sum_{j=0}^{d}a_js^j$ with nonnegative coefficients,
ordinary ultra-log-concavity means that the sequence
\begin{equation}\label{eq:normalization}
 b_j=\frac{a_j}{\binom dj},\qquad 0\le j\le d,
\end{equation}
is log-concave: $b_j^2\ge b_{j-1}b_{j+1}$ at every internal index.
We use the degree-based convention of
\cite[Definition~2.4]{Chen2025}. For interval support this is $\ULC(d)$;
Liggett proves convolution closure \cite{Liggett1997}.
Here $d=m-1$, even when the constant coefficient vanishes. In particular,
normalization by $\binom mj$ would give a different question.

\begin{theorem}\label{thm:main}
Let $m\ge4$ be an integer and let $P$ be either polynomial in
\eqref{eq:families}. Then $P$ is monic of degree $m-1$ with nonnegative
integer coefficients. Writing $a_j=[s^j]P$, one has
\begin{equation}\label{eq:main}
 j(m-1-j)a_j^2<(j+1)(m-j)a_{j-1}a_{j+1}
 \qquad(4\le j\le m-2).
\end{equation}
Moreover, for every $m\ge4$,
\begin{equation}\label{eq:endpoint-statement}
 (m-2)a_{m-2}^2-2(m-1)a_{m-3}a_{m-1}=-m.
\end{equation}
Consequently, neither $f_m$ nor $h_m$ is ordinary ultra-log-concave
for any $m\ge4$.
\end{theorem}

The final non-ULC assertion is already contained in Chen's
Proposition~3.8. The contribution here is the uniform strict inequality
\eqref{eq:main}, posed as Conjecture~3.9 after computations through
$m=100$ \cite[p.~16]{Chen2025}. All these locators refer specifically
to arXiv:2503.16296v2. Theorem~\ref{thm:main} proves that source statement
in its full range. Chen's Proposition~3.10 concerns $\ULC(m)$ and
therefore does not supply these degree-$(m-1)$ inequalities.
A journal version is recorded under
\href{https://doi.org/10.1016/j.geomphys.2025.105666}{doi:10.1016/j.geomphys.2025.105666}.
Its complete text was not inspected in preparing this note; no assertion
of priority relative to that version or all subsequent literature is made.

Moment determinants provide a familiar connection between measures and
log-convexity; Wang and Zhu give a combinatorial treatment, including
infinite log-convexity for Stieltjes moment sequences
\cite[Sections~1--2]{WangZhu2016}. Our reduction involves a density that
changes sign in one parity. The essential step is to transfer positivity
from a different, nonnegative density through an exact determinant identity.

\section{A moment determinant}

For nonnegative integers $p,k,r$, define
\begin{equation}\label{eq:moments}
 M_r(p,k)=\int_{-1}^{1}x^{k+r}(1+x)^p\dd x,
 \qquad H(p,k)=M_0(p,k)M_2(p,k)-M_1(p,k)^2.
\end{equation}

\begin{lemma}\label{lem:positive}
For all integers $p\ge3$ and $k\ge0$, one has $H(p,k)>0$.
\end{lemma}

\begin{proof}
If $k$ is even, $x^k(1+x)^p\dd x$ is a nonnegative measure on $[-1,1]$.
Its density is positive on $(0,1)$, where the functions $1$ and $x$
are not proportional almost everywhere. Strict Cauchy--Schwarz gives
$M_1^2<M_0M_2$.

Fix an odd integer $k\ge1$. For each integer $p\ge1$, all the moments
$M_r(p,k)$ are positive. When $k+r$ is even, the integrand is
nonnegative and positive on $(0,1)$. When $k+r$ is odd, reflection gives
\begin{equation}\label{eq:reflection}
 M_r(p,k)=\int_0^1 x^{k+r}\bigl((1+x)^p-(1-x)^p\bigr)\dd x>0.
\end{equation}
In this parity, introduce the auxiliary determinant
\begin{equation}\label{eq:auxiliary}
 E(p,k)=M_1(p,k)M_3(p,k)-M_2(p,k)^2>0.
\end{equation}
Indeed, the density $x^{k+1}(1+x)^p$ is nonnegative and positive on
$(0,1)$, with moments of orders $0,1,2$ equal to $M_1,M_2,M_3$.
Thus the same strict Cauchy--Schwarz argument proves
\eqref{eq:auxiliary}. It is applied only to this nonnegative measure.

For the base $p=3$, expansion of $(1+x)^3$ yields
\[
 M_0=\frac{6}{k+2}+\frac{2}{k+4},\qquad
 M_1=\frac{2}{k+2}+\frac{6}{k+4},\qquad
 M_2=\frac{6}{k+4}+\frac{2}{k+6}.
\]
Consequently,
\begin{equation}\label{eq:base}
 H(3,k)=\frac{64(k+1)}{(k+2)^2(k+4)^2(k+6)}>0.
\end{equation}
For example, after setting $t=k+2$, the numerator identity needed here is
$t(2t+3)(2t+7)-(t+4)(2t+1)^2=4(t-1)$.

Multiplication of the integrand by $1+x$ gives
\begin{equation}\label{eq:moment-step}
 M_r(p+1,k)=M_r(p,k)+M_{r+1}(p,k).
\end{equation}
Abbreviate the four moments at $(p,k)$ by $A,B,C,D$, in that order,
and put $H=AC-B^2$, $E=BD-C^2$. Since $B>0$, direct algebra gives
\begin{align}
 H(p+1,k)&=(A+B)(C+D)-(B+C)^2\notag\\
 &=H+(AD-BC)+E\notag\\
 &=\left(1+\frac CB\right)H+
   \left(1+\frac AB\right)E.\label{eq:induction}
\end{align}
The last equality follows from $B(AD-BC)=CH+AE$.
Both displayed coefficients are positive, as is $E$ by
\eqref{eq:auxiliary}. Starting from \eqref{eq:base}, induction proves
$H(p,k)>0$ for every $p\ge3$ at this fixed odd $k$. Together with
the even case, this proves the lemma.
\end{proof}

\section{Proof of the coefficient inequalities}

\begin{proof}[Proof of Theorem~\ref{thm:main}]
Geometric summation gives, for $r\ge1$,
\begin{equation}\label{eq:positive-coefficients}
 f_{2r}(s)=s\sum_{\ell=0}^{r-1}(s+1)^{2\ell},\qquad
 f_{2r+1}(s)=1+s(s+1)\sum_{\ell=0}^{r-1}(s+1)^{2\ell}.
\end{equation}
Also, subtraction in \eqref{eq:families} gives
\begin{equation}\label{eq:shared}
 h_m(s)=f_m(s)+(-1)^m.
\end{equation}
These identities establish polynomiality, nonnegative integer
coefficients, degree $m-1$, and monicity for both families. For odd $m$,
\eqref{eq:shared} removes the constant $1$ of $f_m$; for even $m$,
it adds a constant $1$. All positive-degree coefficients agree.
Thus every coefficient difference under consideration is the same for
the two families, since its index is at least $2$.

We first treat the endpoint, which will also cover $m=4,5$.
Comparison of the top three coefficients in
$(s+2)f_m=(s+1)^m-(-1)^m$ gives
\begin{equation}\label{eq:top}
 a_{m-1}=1,\qquad a_{m-2}=m-2,\qquad
 a_{m-3}=\frac{m^2-5m+8}{2}.
\end{equation}
All three indices are positive for $m\ge4$, so \eqref{eq:top} holds
for $h_m$ as well. Define
\begin{equation}\label{eq:difference}
 D_{m,j}=j(m-1-j)a_j^2-(j+1)(m-j)a_{j-1}a_{j+1}.
\end{equation}
Substitution into \eqref{eq:difference} gives
\begin{equation}\label{eq:endpoint}
 D_{m,m-2}=(m-2)^3-(m-1)(m^2-5m+8)=-m.
\end{equation}
This is a failing internal index for every $m\ge4$.

It remains to prove \eqref{eq:main}. Work with $f_m$ and set
$b_j=a_j/\binom{m-1}{j}$. The binomial coefficient ratios give
\begin{equation}\label{eq:minor-normalized}
 D_{m,j}=j(m-1-j)\binom{m-1}{j}^{\!2}
 \bigl(b_j^2-b_{j-1}b_{j+1}\bigr)
\end{equation}
at each internal index. The prefactor is strictly positive.
The fundamental theorem of calculus gives the polynomial identity
\begin{equation}\label{eq:polynomial-integral}
 f_m(s)=m\int_0^1\bigl(-1+(s+2)t\bigr)^{m-1}\dd t.
\end{equation}
For $s\ne-2$ this follows by direct integration; equality of
polynomials includes $s=-2$. Expanding
$-1+(s+2)t=st+(2t-1)$ in the finite binomial sum, then making the
substitution $x=2t-1$, yields
\begin{equation}\label{eq:coefficient-integral}
 b_j=m\int_0^1t^j(2t-1)^{m-1-j}\dd t
 =\frac{m}{2^{j+1}}\int_{-1}^1(1+x)^jx^{m-1-j}\dd x.
\end{equation}

For $4\le j\le m-2$, put $p=j-1$ and $k=m-j-2$.
Then $p\ge3$, $k\ge0$; conversely every such pair corresponds to
$j=p+1$, $m=p+k+3$ in the target range. With $M_r=M_r(p,k)$,
\eqref{eq:coefficient-integral} gives
\begin{equation}\label{eq:adjacent}
 \begin{aligned}
 b_{j-1}&=\frac{m}{2^j}M_2,\\
 b_j&=\frac{m}{2^{j+1}}(M_1+M_2),\\
 b_{j+1}&=\frac{m}{2^{j+2}}(M_0+2M_1+M_2).
 \end{aligned}
\end{equation}
Multiplication and cancellation therefore give
\begin{equation}\label{eq:reduction}
 b_{j-1}b_{j+1}-b_j^2=\frac{m^2}{2^{2j+2}}H(p,k)>0
\end{equation}
by Lemma~\ref{lem:positive}. Equations \eqref{eq:minor-normalized}
and \eqref{eq:reduction} imply $D_{m,j}<0$ throughout the target
range. Equation~\eqref{eq:shared} transfers the inequality to $h_m$,
and \eqref{eq:endpoint} completes the assertion for all $m\ge4$.
\end{proof}

\section{Exact checks and scope}

The higher-index range is empty at $m=4,5$, but
\eqref{eq:endpoint} gives $D_{4,2}=-4$ and $D_{5,3}=-5$.
The lower cutoff cannot be replaced uniformly by $j=3$:
\[
 f_6(s)=3s+6s^2+7s^3+4s^4+s^5,\qquad D_{6,3}=6>0.
\]
Equivalently, the corresponding moment determinant is
$H(2,1)=-16/225$. These calculations illustrate the boundary of the
theorem; they do not alter its quantified range.

Exact sign certificates turn a comparison into rational residuals.
For example, Zeng, Liu and Ratnavelu use seven strict rational
inequalities to certify a Schatten-norm counterexample
\cite[Section~2, equations~(11)--(18)]{ZengLiuRatnavelu2026}.
For the present formulas, moment evaluation itself is rational:
\begin{equation}\label{eq:finite-moment}
 \int_{-1}^1 x^q(1+x)^p\dd x
 =2^{p+1}\sum_{\ell=0}^q
 \frac{\binom q\ell 2^\ell(-1)^{q-\ell}}{p+\ell+1}
 \qquad(p,q\in\mathbb Z_{\ge0}).
\end{equation}
It follows by substituting $x=2t-1$ and expanding $(2t-1)^q$.
Long division of the two original numerators, together with
\eqref{eq:finite-moment}, gives a direct way to check the coefficient
and determinant formulas with integer and rational arithmetic.

A portable Python script, \texttt{verification.py}, is supplied with
the manuscript source. The exact checks cover both families for every
$4\le m\le512$: 1,018 polynomial divisions, 257,556 strict-negative
higher-index differences, and 1,018 endpoint identities. They also
check the coefficient-to-moment reduction at all 703 target pairs for
$6\le m\le42$, and moment positivity at all 920 pairs
$3\le p\le25$, $0\le k\le39$. The latter rectangle includes 460
odd-$k$ induction checks. Separate symbolic checks with SymPy
\cite{Meurer2017} verify the base, endpoint, cancellation and
induction identities. These are finite checks of formulas and indices;
the all-parameter conclusion follows from the proofs above.

The theorem concerns exactly the two families in \eqref{eq:families}
and their degree-based normalization. It asserts neither a general
ultra-log-concavity result for graph classes nor a classification of
all the lower-index signs.

\paragraph{Computational and AI assistance.}
OpenAI Codex assisted with the moment reduction, proof development,
counterexample searches, verification code, literature checks and
manuscript preparation. A separate model context reconstructed the
argument and wrote a second exact verifier. This was an additional
checking pass, not external peer review or statistical independence.
The computations use Python integer and rational arithmetic and SymPy
symbolic algebra. No floating-point test is used to establish a sign,
and no formal proof assistant verification is claimed.

\bibliographystyle{plain}
\bibliography{references}
\end{document}
